\emph{本文针对非线性动态系统状态估计问题，提出一种无需显式求导的新型递推滤波方法——无迹滤波（unscented filter）。传统的线性最小均方递推滤波框架依赖系统动态与观测方程的线性假设；将其推广到非线性系统时，经典的一阶线性化滤波方法在当前估计点附近对非线性函数做一阶截断展开，不仅需要繁琐且易错的一阶偏导数矩阵推导，更会在强非线性条件下因忽略高阶项而产生显著偏差乃至滤波发散。本文基于「近似一个概率分布比近似任意非线性函数更容易」的核心思想，提出无迹变换（unscented transformation）：通过确定性地构造 $2n + 1$ 组对称加权采样点并直接经由真实非线性映射传播，在保持与一阶线性化滤波同阶计算复杂度 $O(n^3)$ 的前提下，对任意非线性变换均能精确匹配至少二阶矩（对对称分布达到三阶精度，对正态分布合理选取缩放参数 $\kappa$ 还可进一步匹配四阶矩）。文中给出了极坐标到直角坐标变换与五维高超声速再入飞行器雷达追踪两个数值算例，验证了无迹滤波在强非线性状态估计中的显著优势。}
